\section{The One Thing It Asks Permission For}
\label{sec:ch17-the-one-thing-it-asks-permission-for}

\index{shareable@\code{"@shareable}!the one refusal in this chapter|(}%
There is a third edit, and it is smaller than either in
section~\ref{sec:ch17-the-field-nobody-approved}. Take the
\code{@override(from: "sessions")} off and leave the declaration standing.
Ratings now says it has an \code{abstract} on \code{Session}, and so does
Sessions, and neither says anything else. Compose:

\begin{minted}[fontsize=\footnotesize,breakanywhere]{text}
The Object "Session" defines the same fields in multiple subgraphs without the "@shareable" directive:
 The field "abstract" is defined in the following subgraphs: "sessions", "ratings".
 However, it is not declared "@shareable" in any of them.
\end{minted}

\index{composition!the shape of the one refusal}%
Exit one. Two services claim one field, the composer has no rule for choosing
between them, and it stops. That is the entire ownership check in Apollo
Federation, and it is worth reading for what it actually says. It says that
two answers to one question is a state the composer cannot plan a route out
of. It has no opinion at all about whether Ratings had any business declaring
the field.

\subsection{Both Sides Have to Say the Word}

\index{shareable@\code{"@shareable}!one document is not enough}%
The obvious repair is to mark the field \code{@shareable}, and the interesting
question is where. Put it on the Sessions copy alone, so that the service that
owns the rows has consented and the newcomer has not, and the composer refuses
again. Its first line is the one above, unchanged; these two are what has
grown a distinction:

\begin{minted}[fontsize=\footnotesize,breakanywhere]{text}
 The field "abstract" is defined and declared "@shareable" in the following subgraph: "sessions".
 However, it is not declared "@shareable" in the following subgraph: "ratings".
\end{minted}

\index{shareable@\code{"@shareable}!the refusal is symmetric}%
Put it on the Ratings copy alone instead and the same message comes back with
the two names exchanged. Both directions refuse, which is the answer to the
question worth asking: there is no asymmetry here to exploit. The service that
owns the rows cannot unilaterally invite a second owner, and a second owner
cannot unilaterally invite itself. Apollo states the rule as a
rule~\autocite{apollovaluetypes}:

\begin{quote}
If a type or field is marked \code{@shareable} in any subgraph, it must be
marked either \code{@shareable} or \code{@external} in every subgraph that
defines it. Otherwise, composition fails.
\end{quote}

\index{composition!enforcing what the vendor documents}%
And this composer does exactly that, which is worth saying out loud after four
chapters of it going its own way. Chapter~\ref{ch:composition} found it
refusing a \code{@cost} weight over a type its own definition disagreed with
Hot Chocolate about, chapter~\ref{ch:second-third-subgraph} found it waving
through a \code{@provides} whose precondition Apollo says is a composition
error, chapter~\ref{ch:null-blast-radius} found it taking the more permissive
of two types for one shareable field, and chapter~\ref{ch:satisfiability}
found it approving a graph with no route through it at all. Here it enforces
the documented rule in both directions and reports which side is missing. Write
both words and the graph composes at exit zero, with both subgraphs recorded
as able to answer the field.

\subsection{What Sharing Actually Promises}

\index{shareable@\code{"@shareable}!nothing checks that the two agree}%
What neither vendor will do for you is the part that matters. Composition
checks that the two declarations have the same shape. Nothing checks, and
nothing can check, that the two resolvers return the same value. WunderGraph
says so in the reference page for the directive~\autocite{cosmoshareable}:

\begin{quote}
However, it is the developer's responsibility to ensure that any such field
definitions always resolve identically, regardless of which subgraph is
ultimately resolving the field. Failure to make each resolver identical would
cause different data to be returned indeterminately.
\end{quote}

\index{shareable@\code{"@shareable}!the promise this book already broke}%
This book has been carrying a field like that since
chapter~\ref{ch:composition}. \code{Query.node} is shareable in two subgraphs,
because chapter~\ref{ch:composition} shipped a type interceptor to make it so,
and
chapter~\ref{ch:second-third-subgraph} then measured that neither subgraph can
keep the promise: Sessions resolves \code{node(id:)} for a session, Speakers
for a speaker, and three requests for the same id disagree with each other
depending on how the client wrote the fragment. \enquote{Indeterminately} is
exactly the word. Chapter~\ref{ch:satisfiability} collected the second bill
for the same declaration, when the shareable root field turned out to buy
silence out of the resolvability check.

\index{ownership!consent to share, none to take}%
So the shape of the whole mechanism is this. To have two owners of a field,
both documents must say so, in a word a human types, and the composer will not
proceed until they agree. To have a different owner of a field, one document
says so and nobody else is consulted at all. Consent is required for the
arrangement where the two teams have to talk anyway, and not required for the
one where a value silently starts coming from somewhere else.
\index{shareable@\code{"@shareable}!the one refusal in this chapter|)}%
